\newcommand{\ModuleID}{EL-II.4}
\newcommand{\ModuleTitle}{Participles and the Ablative Absolute}
\newcommand{\ModuleStage}{Stage II — Core Grammar}
\newcommand{\ModuleUnit}{II.4}
\newcommand{\ModuleOutcome}{Form and attach participles, preserve their relative time and voice, and expand or compress ablative absolutes without changing agency}
\newcommand{\ModulePrerequisites}{EL-II.3 mastery: distinguish infinitive structures and preserve relative time}
\newcommand{\ModuleSessionOne}{Form and decline the ordinary participles and apply the attachment checklist, then complete Worksheet II.13 as the guided application.}
\newcommand{\ModuleSessionTwo}{Expand and compress participial and ablative-absolute pairs while preserving controller, agency, and chronology, then complete Worksheet II.16 independently.}
\newcommand{\ModulePrintedPractice}{Worksheets II.13 (guided Variant A) and II.16 (independent Variant D), with terminal keys}
\newcommand{\ModuleExtensionPractice}{Worksheets II.14 and II.15 remain optional remediation or delayed-recall variants}
\newcommand{\ModuleNext}{EL-II.5, Subjunctive Morphology and Sequence of Tenses}
\newcommand{\ModuleMastery}{Form and decline at least 90 percent of fifteen prompted participles, identify every participle's controller and relative time in an unseen passage, and convert four finite/absolute pairs without changing agency or chronology.}
\newcommand{\ModuleDelayNote}{}
\newcommand{\ModuleReferenceFocus}{%
  \begin{itemize}
    \item present active, perfect passive, and future active participle formation;
    \item active-perfect participles of deponents and the lesson's irregular forms;
    \item agreement, controller, attachment, voice, and relative time; and
    \item ablative-absolute noun/participle coordinates and finite-clause expansion.
  \end{itemize}}
\newcommand{\ModuleContrast}{A participle's agreement identifies its controller; its tense gives time relative to the main verb. An ablative absolute is grammatically detached from the main clause and must not be assigned a false main-clause subject.}
\newcommand{\ModuleLessonSource}{\CourseRoot/shared/content/volume2/all}
\newcommand{\ModuleExerciseSource}{\CourseRoot/shared/exercises/volume2/all}
\newcommand{\ModuleWorkedBank}{II}
\newcommand{\ModuleExerciseBank}{II}
\newcommand{\ModuleWorksheetVariants}{A,D}
\newcommand{\ModuleScopeNote}{This packet teaches the central participial system; contextual English may vary while controller, voice, and chronology remain fixed.}
\newcommand{\ModuleEnrichment}{%
  \SubmoduleExtension
    {Forms and logic}
    {An attachment checklist makes controller, agreement, voice, and relative time explicit before interpretation.}
    {For each participle type in the lesson, fill five boxes: dictionary verb, morphology, agreeing controller, voice, and time relative to the finite verb. Include one deponent perfect participle.}
    {Explain why perfect morphology does not always mean passive sense in a deponent.}
    {Every ordinary perfect passive participle carries passive voice, but a deponent's perfect participle has the active lexical sense identified by its dictionary entry. Agreement and relative time remain independently stated.}
  \SubmoduleExtension
    {Reading the forms in context}
    {A participle can be attributive, circumstantial, supplementary, or periphrastic even when its morphology is unchanged.}
    {Classify the lesson's contextual participles by attachment and use. Draw an arrow to each controller and write the clause evidence that rules out one neighboring classification.}
    {Name the evidence needed before translating a participle as a full subordinate clause.}
    {The answer must identify controller agreement, finite verb, semantic relation, and discourse context. A subordinate-clause English rendering is acceptable only when it preserves the participle's voice and relative time without inventing an unsupported relation.}
  \SubmoduleExtension
    {Writing laboratory}
    {Expansion pairs expose whether an ablative absolute retains its own controller and chronology.}
    {Use the lesson's finite-to-absolute patterns to create one present-participle and one perfect-participle pair. Label the controller and relative time on both sides, then reverse the transformation.}
    {State what must remain unchanged through both directions.}
    {The same proposition, agency, controller, voice, and relative chronology must survive. The absolute pair uses ablative agreement and remains grammatically detached; its finite expansion may choose a contextually justified conjunction but may not invent causality or concession.}
}
